\documentclass[11pt]{article}
\usepackage[T1]{fontenc}
\usepackage{lmodern}
\usepackage{amsmath,amssymb,amsthm}
\usepackage[margin=1in]{geometry}
\usepackage[hidelinks]{hyperref}
\usepackage{microtype}

\newtheorem{theorem}{Theorem}
\newtheorem{lemma}[theorem]{Lemma}
\newcommand{\NN}{\mathbb N_0}
\newcommand{\OK}{\mathcal O_K}

\title{Minimum Shift-Equivalence Lag for the $2\times2$ $k=11$ Baker Pair:\\
A Partial Result on Strong Shift Equivalence}
\author{Xiaojian Zeng\\
PBC School of Finance, Tsinghua University, Beijing 100190, China\\
\texttt{xiaojianzeng@foxmail.com}}
\date{}

\begin{document}
\maketitle

\begin{abstract}
For the Baker pair
\[
A=\begin{pmatrix}1&12\\11&1\end{pmatrix},\qquad
B=\begin{pmatrix}1&132\\1&1\end{pmatrix},
\]
we prove that a shift equivalence over the nonnegative integers whose two
endpoint intertwiners are $2\times2$ has lag at least $22$, and we exhibit such
a shift equivalence at lag $22$.  The proof has no entry bound: it reduces all
possible witnesses to a generalized Pell equation and enumerates its complete
unit orbits.  Since a strong-shift-equivalence chain with $m$ elementary links
composes to a $2\times2$ shift equivalence of lag $m$, every such chain, if one
exists, has at least $22$ links.  Whether $A$ and $B$ are strong shift
equivalent without restrictions on chain length, intermediate dimension, or
entries remains unresolved.
\end{abstract}

\section{Scope and context}

Two nonnegative integral square matrices are elementarily strong shift
equivalent if they are respectively $RS$ and $SR$ for nonnegative integral
rectangular matrices $R,S$; strong shift equivalence (SSE) is the equivalence
relation generated by these moves.  Williams introduced this matrix relation
in the classification of shifts of finite type \cite{Williams1973,Williams1974}.
For square matrices $A,B$, a shift equivalence (SE) of lag $L$ consists of
nonnegative integral matrices $R,S$ satisfying
\begin{equation}\label{eq:se}
 AR=RB,\qquad SA=BS,\qquad RS=A^L,\qquad SR=B^L.
\end{equation}
Baker established early sufficient conditions in the $2\times2$ setting
\cite{Baker1983,Baker1987}, and Choe and Shin later extended part of the
negative-determinant range \cite{ChoeShin1997}.  Those results do not reach
the present determinant $-131$ (which is prime in absolute value).

Jeandel's recent computational study treats the family
\[
A_k=\begin{pmatrix}1&k+1\\k&1\end{pmatrix},\qquad
B_k=\begin{pmatrix}1&k(k+1)\\1&1\end{pmatrix}.
\]
It supplies SSE chains for $k\leq10$ and $k=12$, reports no chain found for
the other tested values $k\geq11$, and explicitly does not turn that search
failure into a nonexistence proof \cite[Problem~3 and Theorem~2]{Jeandel2026}.  Its discussion
of lag $4$ likewise gives only preliminary, unproved evidence for which
parameters admit such an SE.  The theorem below instead fixes $k=11$ and
determines the least lag in the precisely delimited $2\times2$ SE class.  It is
a lower bound, not a decision of unrestricted SSE.  A focused literature
search found no equivalent fixed-instance lag theorem, but no claim of
priority is made.

\begin{theorem}\label{thm:main}
Let
\[
A=\begin{pmatrix}1&12\\11&1\end{pmatrix},\qquad
B=\begin{pmatrix}1&132\\1&1\end{pmatrix}.
\]
If $L\geq1$ and $R,S\in M_2(\NN)$ satisfy \eqref{eq:se}, then $L\geq22$.
Such $R,S$ exist for $L=22$.  Consequently, every SSE chain from $A$ to $B$,
if one exists, has at least $22$ elementary links.
\end{theorem}

\section{Reduction to a norm equation}

Write
\begin{equation}\label{eq:power}
A^L=\begin{pmatrix}a&12b\\11b&a\end{pmatrix},\qquad
a+b\sqrt{132}=(1+\sqrt{132})^L.
\end{equation}
Thus $a^2-132b^2=(-131)^L$.  Comparing entries in the first two identities
of \eqref{eq:se} gives, with no invertibility assumption,
\begin{equation}\label{eq:forms}
R=\begin{pmatrix}x&12z\\z&11x\end{pmatrix},\qquad
S=\begin{pmatrix}11u&12v\\v&u\end{pmatrix}
\end{equation}
for $x,z,u,v\in\NN$.  Put
\[
(p,q,r,s)=(xu,zv,xv,zu).
\]
Direct multiplication shows that both product identities in \eqref{eq:se}
are equivalent to
\begin{equation}\label{eq:scalar}
11p+12q=a,\qquad r+s=b,\qquad pq=rs.
\end{equation}

This scalar reduction is reversible.  Indeed, if a nonnegative integral
quadruple satisfies \eqref{eq:scalar}, then
$\bigl(\begin{smallmatrix}p&r\\s&q\end{smallmatrix}\bigr)$ has determinant
zero and has a nonnegative integral outer factorization.  To see this when
its first row is nonzero, let $g=\gcd(p,r)$ and write
$(p,r)=g(u,v)$ with $\gcd(u,v)=1$.  The equality $pq=rs$ gives
$uq=vs$, hence $(s,q)=z(u,v)$ for some $z\in\NN$; the zero-row cases are
immediate.  Taking $x=g$ recovers \eqref{eq:forms} and all four identities.

For a scalar solution set
\[
X=24q-a,\qquad Y=|r-s|.
\]
Completing the square in \eqref{eq:scalar} yields
\begin{equation}\label{eq:pell}
X^2-132Y^2=(-131)^L.
\end{equation}
Conversely, a solution of \eqref{eq:pell} returns a scalar solution precisely
when
\begin{equation}\label{eq:reconstruct}
q=\frac{X+a}{24},\quad p=\frac{a-12q}{11},\quad
r=\frac{b+Y}{2},\quad s=\frac{b-Y}{2}
\end{equation}
are nonnegative integers.  In particular, every admissible solution has
$|X|\leq a$ and $0\leq Y\leq b$.

\section{Complete unit orbits}

Let $K=\mathbb Q(\sqrt{33})$ and
$\OK=\mathbb Z[(1+\sqrt{33})/2]$, whose discriminant is $33$.  Minkowski's
bound is $\sqrt{33}/2<3$.  Every ideal class therefore has an integral
representative of norm $1$ or $2$.  Since $33\equiv1\pmod 8$, the prime $2$
splits in $K$, and since
\[
N\!\left(\frac{5+\sqrt{33}}2\right)=-2,
\]
the two primes above $2$ are principal.  Hence $\OK$ has class number one.

Set
\[
\lambda=1+2\sqrt{33},\qquad \bar\lambda=1-2\sqrt{33},\qquad
\epsilon=23+4\sqrt{33}.
\]
The conjugate elements $\lambda,\bar\lambda$ have norm $-131$ and generate
the two distinct primes above $131$.  Moreover, $\epsilon$ is the least unit
greater than $1$.  For completeness, a unit $(u+v\sqrt{33})/2$ satisfies
$u^2-33v^2=\pm4$.  For $1\leq v<8$, none of
$33v^2+4$ and $33v^2-4$ is a square; for $v=8$, the plus equation has
$(u,v)=(46,8)$ and the minus equation has no solution.  Larger positive $v$
give larger positive units.  The standard cyclicity theorem for units of a
real quadratic field \cite[Chap.~I]{Neukirch1999} therefore gives all
norm-one units as $\pm\epsilon^n$.

Because the class number is one, unique factorization of ideals now implies
that every integral solution
$X+2Y\sqrt{33}$ of \eqref{eq:pell} has the form
\begin{equation}\label{eq:orbits}
\pm\lambda^j\bar\lambda^{L-j}\epsilon^n,\qquad
0\leq j\leq L,\quad n\in\mathbb Z.
\end{equation}
This is a complete parametrization, not a bounded search in $X$ or $Y$.

The reconstruction bounds make the unit exponent finite.  Both real
embeddings of an admissible $X+2Y\sqrt{33}$ have absolute value at most
\[
a+b\sqrt{132}=(1+\sqrt{132})^L.
\]
Every absolute embedding of the nonunit factor in \eqref{eq:orbits} is
greater than $1$, whereas
$\epsilon=23+2\sqrt{132}>1+\sqrt{132}$.  Thus $n\geq L+1$ violates the
first embedding bound and $n\leq-L-1$ violates the second.  Consequently
\begin{equation}\label{eq:finite}
0\leq j\leq L,\qquad -L\leq n\leq L,
\end{equation}
together with both signs, is exhaustive.

\section{Exact enumeration and the lag-\texorpdfstring{$22$}{22} witness}

The finite enumeration is reproducible with integer arithmetic alone.  A pair
$(c,d)$ represents $c+d\sqrt{132}$ and is multiplied by
\[
(c,d)(e,f)=(ce+132df,\,cf+de).
\]
For each $L$, we compute $(a,b)=(1,1)^L$, loop over the two signs and the
ranges $0\leq j\leq L$ and $-L\leq n\leq L$ in
\eqref{eq:orbits}, retain $|X|\leq a$ and $0\leq Y\leq b$, and apply the
integrality and nonnegativity tests in \eqref{eq:reconstruct}.  The resulting
numbers of scalar solutions for $L=1,\ldots,22$ are
\[
(0,0,0,0,0,0,0,0,0,0,0,0,0,0,0,0,0,0,0,0,0,1).
\]
At lag $22$ the scalar solution is
\begin{align*}
p=s&=330834707029117152443,\\
q=r&=56286015521267519916875.
\end{align*}
It lifts to
\[
R=\begin{pmatrix}1&12\\1&11\end{pmatrix},\qquad
S=\begin{pmatrix}
3639181777320288676873&675432186255210239002500\\
56286015521267519916875&330834707029117152443
\end{pmatrix}.
\]
All entries are nonnegative and direct integer multiplication verifies the
four identities in \eqref{eq:se}.  The complete orbit argument excludes every
smaller lag, proving the first two assertions of Theorem~\ref{thm:main}.

Finally, suppose an SSE chain has $m$ elementary links, with factors
$M_{i-1}=R_iS_i$ and $M_i=S_iR_i$.  The products
$R=R_1\cdots R_m$ and $S=S_m\cdots S_1$ are $2\times2$ at the endpoints and
satisfy \eqref{eq:se} with lag $m$.  Hence $m\geq22$.  Intermediate matrices
may have arbitrary dimensions and entries; the conclusion is solely a lower
bound on chain length.  Existence of an unrestricted chain is not proved or
disproved.

\section*{Computational reproducibility and assistance disclosure}

The accompanying dependency-free verifier evaluates quadratic-ring products
as integer pairs, enumerates the proved finite ranges, checks the reconstruction
filters, and verifies the displayed matrices.  It also directly enumerates the
original scalar parameter through lag $8$ as a separate cross-check.  Thus
floating-point arithmetic and randomization play no role in the certificate.
For comparison, Zeng, Liu, and Ratnavelu use explicit rational inequalities to
certify a different $2\times2$ matrix assertion
\cite{ZengLiuRatnavelu2026}; that preprint is cited only for this
methodological parallel, not as prior work on shift equivalence.

The proof was developed and the manuscript drafted with assistance from an
OpenAI language model.  Exact-arithmetic programs were used to discover and
replay the finite orbit calculation.  A separate language-model context
reconstructed the proof and ran an independently implemented enumeration over
a wider exponent window; this was an internal audit, not peer review.  These
checks support only the explicit calculations; the completeness of the search
rests on the class-number and unit-exponent arguments above.  The author takes
responsibility for every statement.

\begingroup
\small
\bibliographystyle{plainurl}
\bibliography{references}
\endgroup

\end{document}
